\subsubsection{Euler Tour Technique (2 Variantes)}
\begin{lstlisting}[style=mystyle, language=C++]
// TC: O(N) aplanar el arbol, O(log N) por query con SegTree / BIT
// Aplana el arbol para convertir problemas de arboles en consultas de rangos
#include <bits/stdc++.h>
using namespace std;

// ============================================================
// VARIANTE 1: Subtree Queries (Subarbol -> Rango contiguo)
// El subarbol de u corresponde al rango [tin[u], tout[u]] (tamano N)
// ============================================================
vector<int> tin, tout;
int t = -1;

void dfs_subtree(int node, int par, vector<int> adj[]){
  tin[node] = ++t;
  for(auto adjN : adj[node]) if(adjN != par){
    dfs_subtree(adjN, node, adj);
  }
  tout[node] = t; // al salir, 't' actual cubre todo el subarbol
}
// En el SegTree/BIT de tamano N:
//   pos tin[u] guarda el valor a[u]
// Query subarbol de u: query_segtree(tin[u], tout[u])
// Update valor de u:   update_segtree(tin[u], val)

// ============================================================
// VARIANTE 2: Path Queries (Camino Raiz -> Nodo u con + / -)
// tin tiene +val[u] y tout tiene -val[u] (tamano 2N)
// Camino de raiz a u = suma en rango [tin[root], tin[u]]
// ============================================================
void dfs_path(int node, int par, vector<int> adj[]){
  tin[node] = ++t;
  for(auto adjN : adj[node]) if(adjN != par){
    dfs_path(adjN, node, adj);
  }
  tout[node] = ++t; // tin y tout reciben tiempos distintos
}
// En el SegTree/BIT de tamano 2N:
//   pos tin[u]  guarda +a[u]
//   pos tout[u] guarda -a[u]
// Query camino raiz a u: query_segtree(tin[0], tin[u])
// Update valor de u:     sumar delta en tin[u], restar delta en tout[u]
\end{lstlisting}
